\scanpage{111}
% archive foliation: 110

OPS $f\deleted{_{\unsure{\wp}}}$ de la forme $\wp'+\alpha\wp+\beta$
(quitte à multiplier par un scalaire)
\[
(43)\quad\left\{\begin{aligned}
f&=\wp'+\alpha\wp+\beta,\qquad\alpha,\beta\in k\\
\operatorname{div}f&=\{x\}+\{y\}+\{z\}-3\{0_E\}
\end{aligned}\right.
\]
On aura donc
\[
(44)\qquad f(x)=f(y)=f(z)=0\qquad\text{i.e.}
\]
\[
(45)\quad\left\{\begin{aligned}
\wp'(x)+\alpha\wp(x)+\beta&=0\\
\wp'(y)+\alpha\wp(y)+\beta&=0\\
\wp'(z)+\alpha\wp(z)+\beta&=0
\end{aligned}\right.
\]
ce qui signifie aussi que le vecteur colonne
$\begin{pmatrix}\wp'(x)\\\wp'(y)\\\wp'(z)\end{pmatrix}$
s'exprime comme comb. lin. des vecteurs colonnes
$\begin{pmatrix}\wp(x)\\\wp(y)\\\wp(z)\end{pmatrix}$ et
$\begin{pmatrix}1\\1\\1\end{pmatrix}$, \inserted{ce} qui implique que ces trois
vect. colonnes sont lin. dépendants, d'où (42).
Inversement, supposons qu'on ait (42), et $x,y,z$ \deleted{\unsure{égaux}}
\inserted{tous distincts} \deleted{par ex que l'on} \inserted{alors on}
n'a pas $\wp(x)=\wp(y)=\wp(z)$ \deleted{sinon}
\inserted{(sinon $x,y,z$ appartiendraient à une même fibre de $\wp$, donc deux parmi
\unsure{eux} \unsure{seraient égaux})}, \deleted{\unsure{Alors}}
les vect. colonne
$\begin{pmatrix}1\\1\\1\end{pmatrix}\begin{pmatrix}\wp(x)\\\wp(y)\\\wp(z)\end{pmatrix}$
sont lin. indép., et (42) signifie qu'il existe $\alpha,\beta\in k$ tels qu'on ait
(45), i.e. (44), où $f$ est donnée par (43) a).
Comme $x,y,z$ sont distincts, on en conclut bien (43) b), d'où (41).
